\paragraph{Calibration of V/Oct ADCs}
\label{sec:calibration-and-pitch-tracking}

The chosen calibration routine implemented in this project is a simple 2-point calibration, where 2 known input values are converted by the ADC. 
From these values an equation of a straight line can be extracted, which will compensate external factors, like Opamp and electrical inaccuracies, as well as internal ADC offset and gain errors. 

The implemented procedure is described in TI's "General ADC Calibration" article \cite{texasinstrumentsGeneralADCCalibration}.

Two known input voltages \( X_0 \) and \( X_1 \) are applied to the ADC, producing digital codes
\( S_0 \) and \( S_1 \).

Assuming a linear ADC model, the transfer function is
\begin{equation}
S = a X + b
\end{equation}
where \( a \) is the gain (slope) and \( b \) is the offset (intercept at \( X = 0 \)).

The calibration constants are computed as
\begin{equation}
a_c = \frac{S_1 - S_0}{X_1 - X_0}
\end{equation}
\begin{equation}
b_c = S_0 - a_c X_0
\end{equation}

During operation, the calibrated input value is obtained by:
\begin{equation}
X = \frac{S - b_c}{a_c}
\end{equation}

\subparagraph{Translating the calibration to Pitch Tracking}

The general two-point ADC calibration can be directly mapped to the pitch tracking
(V/Oct) system by interpreting the ADC conversion codes as normalized control voltages
and the calibrated output value as pitch measured in semitones.

Two known reference voltages are applied to the V/Oct input:
\begin{itemize}
	\item \( X_0 = 1\,\mathrm{V} \), corresponding to musical note C1
	\item \( X_1 = 3\,\mathrm{V} \), corresponding to musical note C3
\end{itemize}

The corresponding ADC conversion results $c_{1,\mathrm{adc}}$ and $c_{3,\mathrm{adc}}$ are normalised to $[0,1]$ by dividing by $2^{16}$, giving $c_1$ and $c_3$.

Pitch is represented internally in semitones, with
C1 defined as 12 semitones and C3 as 36 semitones.
Assuming a linear relationship between normalized ADC value and pitch, the model becomes:
\begin{equation}
p = m \cdot c + b
\end{equation}
where \( p \) is the pitch in semitones, \( c \) is the normalized ADC value,
\( m \) is the pitch scaling factor, and \( b \) is the pitch offset.

The calibration constants are computed as:
\begin{equation}
m = \frac{36 - 12}{c_3 - c_1} = \frac{24}{c_3 - c_1}
\end{equation}
\begin{equation}
b = 12 - m \cdot c_1
\end{equation}

\paragraph{Pitch Tracking}


During operation, the pitch in semitones is calculated from a measured ADC value \( c \) as:
\begin{equation}
p = m \cdot c + b
\end{equation}

Finally, the pitch is converted to a playback speed factor using the exponential
relationship between pitch and frequency:
\begin{equation}
\text{playback\_speed} = 2^{\frac{p}{12}}
\end{equation}

\paragraph{Limitations}
The pitch offset $b$ is derived by extrapolating the calibration line to 0\,V, making absolute pitch accuracy dependent on the accuracy of the reference voltages applied during calibration. 
Relative pitch tracking between notes is unaffected.

When no cable is connected, the inverting input stage idles at a bias voltage determined by the input resistor network, causing a small pitch offset from true base pitch.
This could be resolved in a future revision by reading the jack detection pin to bypass V/Oct tracking when no cable is inserted.